Niniejszy rozdział szczegółowo opisuje metodę badawczą pracy.
Przedstawia hipotezę badawczą oraz definiuje zmienne kluczowe dla przeprowadzonego eksperymentu.
Na podstawie wstępnego szkicu badania, wskazuje zmienne zakłócające i przedstawia metody ich kontrolowania.
Opisuje kod źródłowy i zadania, które wykonywali badani programiści.
Aby zadbać o rzetelność wyników, rozdział o metodzie badawczej kończy się szczegółowym scenariuszem przeprowadzonych
eksperymentów.

Wyjaśnienia zawarte poniżej mają na celu przede wszystkim spojrzeć na problem ze strony badawczej --- odpowiedzieć
na pytania, dlaczego autor pracy wybrał eksperyment, do jakiego stopnia przewidział zachowanie badanych, w
jaki sposób zaprojektował przeprowadzany eksperyment i jak testuje pozyskane dane pod kątem statystycznej istotności
różnic w nich istniejących.
Nadrzędną wartością tego rozdziału jest odtwarzalność toku myślowego prowadzącego do podjęcia decyzji o
przeprowadzeniu dokładnie tego eksperymentu, umożliwiając jego najdokładniejszą ocenę.

Przedmiotem badań całej pracy jest obciążenie kognitywne programisty podczas wprowadzania zmian w
istniejącym kodzie źródłowym w języku \texttt{Java}.
Wybór \texttt{Java}
jako języka programowania jest nieprzypadkowy --- w Sii sp. z o.o., z którego pochodzą badani, 40\%
stanowisk jest związanych z tworzeniem oprogramowania właśnie w tym języku.
\texttt{Java} jest tam najczęściej używana.
Globalnie zajmuje trzecie miejsce z wynikiem 10,15\% w rankingu najbardziej popularnych języków
programowania~\cite{TiobeIndex2024}, co zwiększa szansę na replikacje badania przedstawionego w niniejszej pracy.

Wspomniane obciążenie kognitywne jest mierzone pośrednio poprzez czas wymagany na zrozumienie kodu i jego modyfikacje.
Badanie koncentruje się na tym, czy wyższy stopień sprzężenia klas wydłuża czas potrzebny programiście na zmianę
kodu według przekazanych wymagań, a to oznaczałoby, że w takiej sytuacji taki kod jest trudniejszy do zrozumienia,
prowadząc do wniosku, że ma to wpływ na obciążenie kognitywne programisty podczas pracy z tak zaprojektowanymi klasami.
Celem badawczym tej pracy jest zweryfikowanie zależności pomiędzy stopniem sprzężenia istniejących klas
a czasem programisty poświęconym na wprowadzenie zmiany.

Autor pracy zaprojektował badanie, korzystając z informacji zawartych w monografii~\cite{wohlin2012experimentation},
podręcznikowi w zakresie badań empirycznych w dziedzinie inżynierii oprogramowania.
